
\textbf{Lemme.} \uline{Soit} @@M0@@ \uline{un morphisme projectif et parfait}, \uline{avec} @@M1@@ \uline{quasi-compact}, \uline{et soit} @@M2@@ \uline{un complexe parfait sur} @@M3@@, \uline{acyclique en degrés} @@M4@@. \uline{Alors il existe un entier} @@M5@@ \uline{tel que pour} @@M6@@ \uline{le complexe} @@M7@@ \uline{soit acyclique en degrés} @@M8@@.

On applique ce lemme au complexe réduit à @@M9@@; il existe donc un entier @@M10@@ tel que pour @@M11@@, on ait @@M12@@ si @@M13@@. On redescend alors par récurrence jusqu'à @@M14@@. Pour cela, on considère la suite exacte:
@@M15@@
et on lui applique le foncteur @@M16@@. Or on a @@M17@@. Comme @@M18@@, @@M19@@ est @@M20@@-acyclique pour @@M21@@. On a donc pour tout @@M22@@, et tout @@M23@@:
@@M24@@
et par conséquent les suites exactes:
@@M25@@
ce qui permet d'achever la récurrence.

ii) Il s'agit de montrer que le morphisme canonique:
@@M26@@
est un isomorphisme. La propriété étant locale, on peut supposer @@M27@@ défini par une suite @@M28@@ de sections de @@M29@@, telles que l'homomorphisme @@M30@@ correspondant soit régulier. Posons: @@M31@@; on définit un homomorphisme d'algèbres: @@M32@@, surjectif, en envoyant 
